\subsection{The obstruction has to survive at the level of queries}
\label{threshold:obstruction-transfer}

Write $\mathcal A(\mu)$ for the following property of a family of input
laws: every $T$-query quantum acceptance function $Q_n:\{-1,1\}^n\to[0,1]$
has a classical decision-tree approximation $D_n$ satisfying
$\mathbb E_{\mu_n}|D_n-Q_n|\le\epsilon$ with at most
$\operatorname{poly}(T,1/\epsilon)$ queries.  The polynomial must be
independent of $n$, but can depend on fixed model parameters.  Decision
trees may be computationally unbounded.  Outputs can be clipped to
$[0,1]$.  Mean-error and the usual high-probability formulation imply
one another with polynomial changes in the accuracy parameters.

\begin{lemma}[Hard-component and perturbation transfer]
\label{threshold:hard-component}
Suppose every classical tree of depth at most $t$ satisfies
$\mathbb E_{\nu_n}|D-Q_n|\ge L_n(t)$.  If
$\mu_n\ge w_n\nu_n$ pointwise, then
\[
 \mathbb E_{\mu_n}|D-Q_n|\ge w_nL_n(t).
\]
If instead $\|\mu_n-\nu_n\|_{\rm TV}\le\eta_n$, then
\[
 \mathbb E_{\mu_n}|D-Q_n|\ge L_n(t)-\eta_n.
\]
The statements also hold for randomized trees, with expectation over
their randomness, whenever the premise does.
\end{lemma}

\begin{proof}
The loss is nonnegative and at most one, so the first assertion follows
by integration and the second by the variational characterization of
total variation.  Conditioning on the random seed proves the randomized
version.
\end{proof}

This identifies two separate requirements for a thermally generated
obstruction.  A hard conditional law must possess an actual quantum
versus classical query separation.  Its Gibbs mass must also be large
enough relative to the separation and the allowed approximation error.
Macroscopic phase coexistence supplies neither requirement automatically.
For example, the exchangeable-simulation theorem in
Section~\ref{exch:heading-2} handles all Curie--Weiss temperatures,
including its two-phase regime.
